% ---- score2 standard preamble (from: python3 score2_tools.py preamble) ----
\documentclass[10pt]{extarticle}
\usepackage[a4paper,margin=1.5cm]{geometry}
\usepackage[T1]{fontenc}
\usepackage{amsmath,amssymb}
\usepackage{graphicx}
\usepackage[export]{adjustbox}
\usepackage{multicol}
\usepackage{enumitem}
\usepackage{booktabs,array,tabularx}
\setlength{\columnsep}{0.7cm}
\setlength{\parindent}{0pt}
\setlength{\parskip}{3pt}
\pagestyle{empty}
\makeatletter
% figure registry (figdefs.tex is generated by the extract step)
\newcommand{\qfigdef}[4]{\@namedef{qf@file@#1}{#2}\@namedef{qf@w@#1}{#3}\@namedef{qf@wide@#1}{#4}}
\newcommand{\qf@check}[1]{\@ifundefined{qf@file@#1}{\PackageError{score2}{Unknown figure id #1}{}}{}}
\newcommand{\qf@img}[4][]{\includegraphics[width=#2\dimexpr\@nameuse{qf@w@#4}pt\relax,max width=#3,max height=0.45\textheight,#1]{\@nameuse{qf@file@#4}}}
% Full-width questions: finalize wraps every question that contains a wide figure in \qfullbegin...\qfullend,
% so the WHOLE question (stem, figures, options) is set across the page in reading order, then the two columns
% resume. The question's head - its number, stem and first figure (\qfullsplit marks the end of the head) -
% is kept on one page, so a number is never left at a page foot with its figure overleaf; the rest may break
% across pages normally, so no big gaps are left. Without a \qfullsplit, the whole question is kept together
% if it is under 45% of a page. Never write these yourself.
\newif\ifqf@full
\newif\ifqf@boxed
\newsavebox{\qf@box}
\newcommand{\qf@open}{\global\qf@fulltrue\par\global\qf@boxedtrue
  \setbox\qf@box\vbox\bgroup\hsize\textwidth\linewidth\textwidth}
\newcommand{\qf@emit}[1]{\par\egroup\global\qf@boxedfalse
  \ifdim\dimexpr\ht\qf@box+\dp\qf@box\relax<#1\textheight\noindent\usebox{\qf@box}\par\else\unvbox\qf@box\fi}
\newcommand{\qfullsplit}{\ifqf@boxed\qf@emit{0.85}\fi}
\newcommand{\qf@close}{\ifqf@boxed\qf@emit{0.45}\fi\global\qf@fullfalse}
\newcommand{\qfullbegin}{\end{multicols}\qf@open}
\newcommand{\qfullend}{\qf@close\begin{multicols}{2}}
\newcommand{\qfullstart}{\qf@open}
\newcommand{\qfulldone}{\qf@close}
% \qfig[scale]{id}: block figure, centred, directly where it is written (size set by the tool).
% Fallback if a wide figure is not inside a full-width question: break out of the columns in place - no floats.
\newcommand{\qfig}[2][1]{\qf@check{#2}%
  \ifqf@full\qf@blk{#1}{#2}\else
  \ifnum\@nameuse{qf@wide@#2}=1 \expandafter\@firstoftwo\else\expandafter\@secondoftwo\fi
  {\end{multicols}\begin{center}\qf@img{#1}{\textwidth}{#2}\end{center}\begin{multicols}{2}}%
  {\qf@blk{#1}{#2}}\fi}
\newcommand{\qf@blk}[2]{\par\smallskip{\centering\qf@img{#1}{\linewidth}{#2}\par}\smallskip}
% \qfigi[scale]{id}: inline figure (inside option lists / tables / minipages); never breaks columns.
\newcommand{\qfigi}[2][1]{\qf@check{#2}\qf@img[valign=c]{#1}{\linewidth}{#2}}
\makeatother
% structure macros - lint reads these to check completeness and figure placement
\newcommand{\papersec}[1]{\par\bigskip{\centering\Large\bfseries PAPER-#1\par}\medskip}
\newcommand{\partsec}[2]{\par\medskip\noindent{\large\bfseries PART-#1: #2}\par\smallskip}
\newcommand{\sect}[1]{\par\smallskip\noindent\textbf{#1}\par}
\newcommand{\qn}[1]{\par\smallskip\noindent\textbf{#1.}~}
\newcommand{\paraq}[2]{\par\smallskip\noindent\textbf{Paragraph for Questions #1 and #2}\par}
\newenvironment{qoptions}{\begin{itemize}[leftmargin=1.8em,itemsep=1pt,topsep=2pt,label={}]}{\end{itemize}}
\newcommand{\opt}[2]{\item[(#1)] #2}
% figure registry (generated by score2_tools.py finalize)
\qfigdef{p1_phy_s1_q04_1}{Score-II - Test 23 (17 Mar 2024)-img/p1_phy_s1_q04_1.png}{122.2}{0}
\qfigdef{p1_phy_s1_q05_1}{Score-II - Test 23 (17 Mar 2024)-img/p1_phy_s1_q05_1.png}{160.1}{0}
\qfigdef{p1_phy_s1_q06_1}{Score-II - Test 23 (17 Mar 2024)-img/p1_phy_s1_q06_1.png}{113.7}{0}
\qfigdef{p1_phy_s2_q02_1}{Score-II - Test 23 (17 Mar 2024)-img/p1_phy_s2_q02_1.png}{105.8}{0}
\qfigdef{p1_phy_s2_q03_1}{Score-II - Test 23 (17 Mar 2024)-img/p1_phy_s2_q03_1.png}{160.4}{0}
\qfigdef{p1_phy_s2_q05_1}{Score-II - Test 23 (17 Mar 2024)-img/p1_phy_s2_q05_1.png}{109.4}{0}
\qfigdef{p1_phy_s3_q01_1}{Score-II - Test 23 (17 Mar 2024)-img/p1_phy_s3_q01_1.png}{96.8}{0}
\qfigdef{p1_phy_s3_q02_1}{Score-II - Test 23 (17 Mar 2024)-img/p1_phy_s3_q02_1.png}{86.9}{0}
\qfigdef{p1_phy_s3_q03_1}{Score-II - Test 23 (17 Mar 2024)-img/p1_phy_s3_q03_1.png}{191.3}{0}
\qfigdef{p1_phy_s3_q04_1}{Score-II - Test 23 (17 Mar 2024)-img/p1_phy_s3_q04_1.png}{122.2}{0}
\qfigdef{p1_phy_s3_q05_1}{Score-II - Test 23 (17 Mar 2024)-img/p1_phy_s3_q05_1.png}{137.4}{0}
\qfigdef{p1_phy_s3_q06_1}{Score-II - Test 23 (17 Mar 2024)-img/p1_phy_s3_q06_1.png}{159.7}{0}
\qfigdef{p1_chem_s1_q03_A}{Score-II - Test 23 (17 Mar 2024)-img/p1_chem_s1_q03_A.png}{86.7}{0}
\qfigdef{p1_chem_s1_q03_B}{Score-II - Test 23 (17 Mar 2024)-img/p1_chem_s1_q03_B.png}{81.8}{0}
\qfigdef{p1_chem_s1_q03_C}{Score-II - Test 23 (17 Mar 2024)-img/p1_chem_s1_q03_C.png}{92.3}{0}
\qfigdef{p1_chem_s1_q03_D}{Score-II - Test 23 (17 Mar 2024)-img/p1_chem_s1_q03_D.png}{110.0}{0}
\qfigdef{p1_chem_s1_q05_1}{Score-II - Test 23 (17 Mar 2024)-img/p1_chem_s1_q05_1.png}{240.4}{0}
\qfigdef{p1_chem_s1_q06_1}{Score-II - Test 23 (17 Mar 2024)-img/p1_chem_s1_q06_1.png}{226.1}{0}
\qfigdef{p1_chem_s1_q06_2}{Score-II - Test 23 (17 Mar 2024)-img/p1_chem_s1_q06_2.png}{44.9}{0}
\qfigdef{p1_chem_s1_q06_3}{Score-II - Test 23 (17 Mar 2024)-img/p1_chem_s1_q06_3.png}{45.1}{0}
\qfigdef{p1_chem_s1_q06_4}{Score-II - Test 23 (17 Mar 2024)-img/p1_chem_s1_q06_4.png}{44.9}{0}
\qfigdef{p1_chem_s1_q06_5}{Score-II - Test 23 (17 Mar 2024)-img/p1_chem_s1_q06_5.png}{66.3}{0}
\qfigdef{p1_chem_s2_q01_1}{Score-II - Test 23 (17 Mar 2024)-img/p1_chem_s2_q01_1.png}{130.5}{0}
\qfigdef{p1_chem_s2_q01_2}{Score-II - Test 23 (17 Mar 2024)-img/p1_chem_s2_q01_2.png}{116.6}{0}
\qfigdef{p1_chem_s2_q04_1}{Score-II - Test 23 (17 Mar 2024)-img/p1_chem_s2_q04_1.png}{229.0}{0}
\qfigdef{p1_chem_s3_q01_1}{Score-II - Test 23 (17 Mar 2024)-img/p1_chem_s3_q01_1.png}{119.9}{0}
\qfigdef{p1_chem_s3_q02_1}{Score-II - Test 23 (17 Mar 2024)-img/p1_chem_s3_q02_1.png}{90.4}{0}
% ---- body skeleton ----
% \begin{document}
% \begin{center}{\Large\bfseries JEE (Main+Advanced) Score-II --- Test NN (DD Mon YYYY)}\end{center}
% \begin{multicols}{2}
%   \papersec{1}   (only for two-paper tests)
%   \partsec{1}{PHYSICS} \sect{Section-I (i): Single Correct} \qn{1} ... \qfig{p1_phy_q01_1} ...
% \end{multicols}
% \end{document}

\begin{document}
\begin{center}{\Large\bfseries JEE (Main+Advanced) Score-II --- Test 23 (17 Mar 2024)}\end{center}
\begin{multicols}{2}

\partsec{1}{PHYSICS}
\sect{Section-I: One or More Correct}

\qn{1} An coil having inductance L and resistance R is in series with a capacitor and an AC source. At resonance, the rms voltage across the inductor is V$_1$ and across the capacitor is V$_2$.
\begin{qoptions}
\opt{A}{The rms current through the circuit is $\dfrac{\sqrt{V_1^2-V_2^2}}{R}$.}
\opt{B}{The rms current through the circuit is $\dfrac{V_1-V_2}{R}$.}
\opt{C}{The quality factor of circuit is $\dfrac{V_2}{\sqrt{V_1^2-V_2^2}}$.}
\opt{D}{The power loss is circuit at resonance is $\dfrac{(V_1-V_2)^2}{R}$.}
\end{qoptions}

\qn{2} The electric field due to a plane electromagnetic wave in air is
\[\vec E = \left[4\cos(10^7\pi t - kz)\hat\imath + 4\sin(10^7\pi t - kz)\hat\jmath\right]\ V/m\]
\begin{qoptions}
\opt{A}{The x component of magnetic field is $B_x = \dfrac43\times10^{-8}\sin(10^7\pi t-kz)T\,\hat\imath$}
\opt{B}{The y component of magnetic field is $B_y = \dfrac43\times10^{-8}\cos(10^7\pi t-kz)T\,\hat\jmath$}
\opt{C}{The poynting vector is $\dfrac{2}{15\pi}\hat k\ W/m^2$.}
\opt{D}{The power transmitted through 1 $m^2$ of a plane area having intercepts $1:2:3$ on the x-axis, y-axis and z-axis respectively is $\dfrac{4}{105\pi}$ W.}
\end{qoptions}

\qn{3} A heat engine operates between two reservoirs at temperatures of 500 K and 40 K. The engine drives a refrigerator which operates between reservoirs at temperatures of 40 K and 20 K. The heat transfer to the heat engine is 2000 kJ and the net work output of the combined engine refrigerator plant is 340 kJ.
\begin{qoptions}
\opt{A}{If both the engine and refrigerator follow the carnot cycle, then the heat transferred to the reservoir at 40 K is 3160 kJ.}
\opt{B}{If both the engine and refrigerator follow the carnot cycle, then the heat transferred from reservoir at 20 K is 1500 kJ.}
\opt{C}{If the efficiency of heat engine and COP of refrigerator both are 50\% of the maximum possible values, the heat transferred to the reservoir at 40 K is 1950 kJ.}
\opt{D}{If the efficiency of heat engine and COP of refrigerator both are 50\% of the maximum possible values, the heat transferred from reservoir at 20 K is 290 kJ.}
\end{qoptions}

\qn{4} As shown in figure the candle is at the center of curvature of the concave mirror, whose focal length is 10.0 cm. The converging lens has a focal length of 32.0 cm and is 85.0 cm to the right of the candle. The candle is viewed looking through the lens from the right. The lens forms two images of the candle. The first is formed by light passing directly through the lens. The second image is formed from the light that goes from the candle to the mirror, is reflected, and then passes through the lens.
\qfig{p1_phy_s1_q04_1}
\begin{qoptions}
\opt{A}{Both images are formed at the same distance from the lens.}
\opt{B}{Both images are real.}
\opt{C}{Both images are inverted with respect to the object.}
\opt{D}{Both images are of same size.}
\end{qoptions}

\qn{5} Figure indicates the lowest energy levels (in electron volts) for five situations in which an electron is trapped in a one-dimensional infinite potential well. In wells B, C, D, and E, the electron is in the ground state. We shall excite the electron in well A to the fourth excited state (at 25 eV). The electron can then de-excite to the ground state by emitting one or more photons, corresponding to one long jump or several short jumps. The photon emitted from well A can be absorbed by electrons in other wells.
\qfig{p1_phy_s1_q05_1}
\begin{qoptions}
\opt{A}{The electron in B will stay in ground state.}
\opt{B}{The electron in C can jump to first excited state.}
\opt{C}{The electron in D can jump to first excited state but not to second excited state.}
\opt{D}{The electron in E can jump to first excited state.}
\end{qoptions}

\qn{6} A step ladder is placed to the smooth horizontal surface of ice. The cross-section of the ladder is the same everywhere and the angle between its legs is $\alpha=60^\circ$. The hinge which connects the two legs is frictionless. The thin, initially tight safety rope breaks.
\qfig{p1_phy_s1_q06_1}
\begin{qoptions}
\opt{A}{The angle between the legs of the ladder when the hinge does not exert any force is $120^\circ$.}
\opt{B}{The force exerted by ice on the ladder is $\dfrac{mg}{2}$.}
\opt{C}{The speed at which the hinge strikes the ice is $\sqrt{\dfrac{3g\ell}{2}}$.}
\opt{D}{The hinge force just before the hinge strikes the ice is $\dfrac{3\sqrt3 mg}{4}$.}
\end{qoptions}

\sect{Section-II: Numerical Value}

\qn{1} A current I flows through a wire of length L shaped as an ellipse of area A. The magnetic field at the centre of ellipse is $\alpha\times10^{-7}$ T. The value of $\alpha$ is. (Take : I = 1A, L = 2m, A = $2m^2$)

\qn{2} Figure shows two circular regions R$_1$ and R$_2$ with radii $r_1=20.0$ cm and $r_2=30.0$ cm. In R$_1$ there is a uniform magnetic field of magnitude B$_1=50.0$ mT directed into the page, and in R$_2$ there is a uniform magnetic field of magnitude B$_2=75.0$ mT directed out of the page (ignore fringing). Both fields are decreasing at the rate of 8.50 T/s. Calculate $\oint\vec E\cdot d\vec s$ for path 1 is $\alpha$V, for path 2 is $\beta$V and path 3 is $\gamma$V, the value of $(\alpha-\beta+\gamma)$ is :- (Take : $\pi=3.14$)
\qfig{p1_phy_s2_q02_1}

\qn{3} The small block at A has a mass of 4 kg and is mounted on the bent rod having negligible mass. If the rotor at B causes a harmonic movement $\delta_B = (0.112\cos15t)$ m, where t is in seconds, determine the steady-state (in cm) amplitude of vibration of the block.
\qfig{p1_phy_s2_q03_1}

\qn{4} A 1.0 $\mu$Bq $^{137}_{55}$Cs source is used for 2.0 hours by a student in a physics lab. Radioactive $^{137}_{55}$Cs decays by $\beta^-$ decay with a half-life of 30 years The average energy of the emitted betas is about 190 keV per decay. The $\beta$ decay is quickly followed by a $\gamma$ wilh an energy of 660 keV. Assuming the student absorbs all emitted energy, the energy absorbed by the student is $x\times10^{-16}$ J. The value of x is.

\qn{5} In a shipping company distribution center, an open cart of mass 50.0 kg is rolling to the left at a speed of 5.00 m/s (Figure). Ignore friction between the cart and the floor. A 15.0-kg package slides down a chute that is inclined at $37^\circ$ from the horizontal and leaves the end of the chute with a speed of 3.00 m/s. The package lands in the cart and they move together. If the lower end of the chute is a vertical distance of 4.00 m above the bottom of the cart, what is the final speed (in m/s) of the cart?
\qfig{p1_phy_s2_q05_1}

\qn{6} Consider a partially polarized light beam, containing a mix of unpolarized and linearly polarized light. The intensity of the beam is analyzed using a linear polarized. At a particular orientation of the polarizer, the outgoing beam has maximum intensity I$_{max}$. Turning the polarizer by a $30^\circ$ angle reduce the outgoing beam's intensity by 10\%. Find the degree of polarization of the beam. $V = \dfrac{I_{max}-I_{min}}{I_{max}+I_{min}}$, where I$_{min}$ is the minimum intensity for any orientation of the polarizer.

\sect{Section-III: Single Digit Integer (0-9)}

\qn{1} Five identical conducting plates each of area A are placed parallel and equidistant at a distance d from each other and are connected as shown. Find the ratio of charges on plates 3 to the charges on plate 5.
\qfig{p1_phy_s3_q01_1}

\qn{2} The current i$_1$ (in A) is :-
\qfig{p1_phy_s3_q02_1}

\qn{3} Figure-(a) shows the interference of two waves having beat frequency $f_1$ and figure-(b) shown the interference of two waves having beat frequency $f_2$. The ratio $\dfrac{f_1}{f_2}$ is :-
\qfig{p1_phy_s3_q03_1}

\qn{4} A gate supporting water is shown in figure. Find the height h of the water so that the gate tips about the hinge. Take the width of the gate as unity.
\qfig{p1_phy_s3_q04_1}

\qn{5} A water drop will last about 1 s on a hot skillet with a temperature between $100^\circ$C and about $200^\circ$C However, if the skillet is much hotter, the drop can last several minutes. The longer lifetime is due to the support of a thin layer of air and water vapor that separates the drop from the metal (by distance L in Fig.). Let L = 0.100 mm, and assume that the drop is flat with height h = 1.50 mm and bottom face area A $=4.00\times10^{-6}$ m$^2$. Also assume that the skillet has a constant temperature T$_s=300^\circ$C and the drop has a temperature of $100^\circ$C. Water has density $\rho=1000\ kg/m^3$, and latent heat of vapourization is $2.2\times10^6$ J/kg, the supporting layer has thermal conductivity k = 0.0275 W/m-K. If conduction is the primary way energy moves from the skillet to the drop, how long (in minutes) will the drop last?
\qfig{p1_phy_s3_q05_1}

\qn{6} When a wood shelf of mass 5.0 kg is fastened inside a slot in a vertical support as shown in figure, the support exerts a torque on the shelf. The torque exerted by the support (about the left end of the shelf) (in N-m) is :-
\qfig{p1_phy_s3_q06_1}

\partsec{2}{CHEMISTRY}
\sect{Section-I: One or More Correct}

\qn{1} SO$_3$ gas is entering the environment at a constant rate of $6.93\times10^{-6}$ g/L/day due to emission of polluting gases from thermal power plant at Kota but at the same time it is decomposing and following first-order kinetics with half-life of 100 days. Based on these details, select the correct statement(s) from the following.
\begin{qoptions}
\opt{A}{Concentration of SO$_3$ in Kota is $1.25\times10^{-5}$ M (assume SO$_3$ present in air reaches steady state).}
\opt{B}{If $10^3$ L of air is passed through 1 L pure water (assuming all SO$_3$ is dissolved in it) and the resulting solution is titrated against 1 N NaOH solution, 15 ml of NaOH solution is required to reach the end point.}
\opt{C}{An industry is manufacturing H$_2$SO$_4$ at the rate of 980 kg per day with the use of SO$_3$ in air and it should use $8\times10^5$ L air per day.}
\opt{D}{If SO$_3$ emission is stopped, then after 1000 days, its concentration will reduce to $1.22\times10^{-8}$ M.}
\end{qoptions}

\qn{2} Decomposition of A$_2$(g) and B$_3$(g) follows first-order kinetics as follows.
\[A_2(g)\to 2A(g);\ K_1(h^{-1}) = e^{-\frac{14000(J)}{RT}+5}\]
\[B_3(g)\to 3B(g);\ K_2(h^{-1}) = e^{-\frac{20000(J)}{RT}+10}\]
Here, K$_1$ and K$_2$ are the rate constants with respect to disappearance of A$_2$ and B$_3$, respectively. One mole, each of A$_2$(g) and B$_3$(g), is taken in a 100 L evacuated flask and at some temperature at which they start decomposing at the same rate. The incorrect information regarding the reactions is
\begin{qoptions}
\opt{A}{The temperature at which the reactions are performed is $\dfrac{1200}{8.314}$K.}
\opt{B}{At any instant, $\left(\dfrac{P_{A_2}}{P_{B_3}}\right)$ will be constant and equal to 1.0, where P$_{A_2}$ and P$_{B_3}$ are the partial pressures of gases at that time.}
\opt{C}{At some instant, the total pressure of all gases in the flask may be less than 0.2 atm.}
\opt{D}{At any instant, $\left(\dfrac{P_A}{P_B}\right)$ will be constant and equal to 1.0, where P$_A$ and P$_B$ are the partial pressures of gases at that time.}
\end{qoptions}

\qn{3} Which of the following reactions correctly represents the change in ammonium salts.
\begin{qoptions}
\opt{A}{\qfigi{p1_chem_s1_q03_A}}
\opt{B}{\qfigi{p1_chem_s1_q03_B}}
\opt{C}{\qfigi{p1_chem_s1_q03_C}}
\opt{D}{\qfigi{p1_chem_s1_q03_D}}
\end{qoptions}

\qn{4} Which of the following statement(s) is/are correct?
\begin{qoptions}
\opt{A}{Sodium carbonate decahydrate effloresces at room temperature to form monohydrate.}
\opt{B}{Lithium can not form ethynide upon reaction with ethyne.}
\opt{C}{Na$^+$ ions are responsible for the transmission of nerve signal.}
\opt{D}{The diluted solution of 'K' in liquid NH$_3$ is blue color due to ammoniated electrons.}
\end{qoptions}

\qn{5} For the given reaction :
\qfig{p1_chem_s1_q05_1}
Correct statement(s) among the following is/are :
\begin{qoptions}
\opt{A}{Product 'R' gives violet colour with FeCl$_3$ test.}
\opt{B}{Product 'Q' gives phenol on reaction with steam}
\opt{C}{Product 'R' gives bromine water test.}
\opt{D}{Product 'T' gives ethyne on reaction with Ag/$\Delta$.}
\end{qoptions}

\qn{6} Consider the following reaction sequence :
\qfig{p1_chem_s1_q06_1}
Choose the correct option :
\begin{qoptions}
\opt{A}{Product 'X' is\\ \qfigi{p1_chem_s1_q06_2}}
\opt{B}{Product 'Y' is\\ \qfigi{p1_chem_s1_q06_3}}
\opt{C}{Product 'Z' is\\ \qfigi{p1_chem_s1_q06_4}}
\opt{D}{Product 'W' is\\ \qfigi{p1_chem_s1_q06_5}}
\end{qoptions}

\sect{Section-II: Numerical Value}

\qn{1} Consider a two chamber vessel. An SPM Separates the two chambers into equal volume as shown
\qfig{p1_chem_s2_q01_1}
A one liter aqueous solution $\mathbf{X}$, (density = 1.6 g/ml) of a non volatile solute is filled in one of the chambers. Soon after that an equal volume of pure water (density = 1.0 g/ml) is filled in the other chamber. Very quickly, the solution surface is now covered with a capillary of 50 $cm^2$ cross section. Pure water now starts flowing from the solvent towards the solution. At equilibrium condition the level of the solution in the capillary is 30.4cm then find the osmotic pressure (in milli-atm) of the solution X. Assume $\Delta V_{mix}=0$

Use: density of Hg = 13.4 g/ml, atmospheric pressure = 76 cm of Hg and $\dfrac{23.6}{176}=0.134$)
(round off to the nearest integer)
\qfig{p1_chem_s2_q01_2}

\qn{2} A sample of Fe$_2$(SO$_4$)$_3$ and FeC$_2$O$_4$ was dissolved in dil. H$_2$SO$_4$. The complete oxidation of reaction mixture required 40 mL of N/16 KMnO$_4$. After the oxidation, the reaction mixture was reduced by Zn and dil. H$_2$SO$_4$. On again oxidation by same KMnO$_4$, 60 mL were required. If the ratio of moles of Fe$_2$(SO$_4$)$_3$ and FeC$_2$O$_4$ in mixture is x : y then find the value of x/y

\qn{3} Orthophosphoric acid on treatment with a reagent (R) in presence of conc. HNO$_3$ gives Canary yellow ppt. (P). The number of transition metal atoms in (P) are x and number of hydrogen atoms in (R) are y. Then find the value of $\dfrac{x}{y}$.

\qn{4} In the following reaction sequence in aq. solution.
\qfig{p1_chem_s2_q04_1}
(i) Number of ligands in A is x.

(ii) G forms a complex with KCN in the presence of air, number of isomers of the complex formed are y

(iii) E on reaction with sodium Nitroprusside forms a complex Na$_\alpha$[Fe(CN)$_\beta$NOS]

Value of $\dfrac{\alpha+\beta+x}{y}$ is.

\qn{5} An organic compound A(C$_2$H$_6$O) on oxidation by PCC give B(C$_2$H$_4$O), which on treatment with NaOH/$\Delta$ give product C with molecular mass 'X'. B on oxidation by KMnO$_4$/H$^+$ produce monobasic carboxylic acid(D) with molecular mass 'Y' with monomeric form. Find the value of X$-$Y.

\qn{6} From the given reaction sequence.
Propyne $\xrightarrow[\text{(ii) CH}_3\text{Br}]{\text{(i) Na}}$ 'X'. If total number of carbon atom in 'X' is A and degree of unsaturation in 'X' is B, then find out (A$-$B).

\sect{Section-III: Single Digit Integer (0-9)}

\qn{1} Quinhydrone is a sparingly soluble one-to-one addition compound formed from hydroquinone and quinone. When solid quinhydrone is dissolved in an aqueous medium, equal concentrations of hydroquinone and quinone result. If a platinum wire is dipped into the solution, the potential of the electrode is governed by the reversible reaction :
\qfig{p1_chem_s3_q01_1}
Where Q denotes quinone and H$_2$Q denotes hydroquinone. The E$^\circ$ value of this half cell reaction is +0.46V with respect to the saturated calomel reference electrode (SCE). Because the potential of the quinine-hydroquinone half reaction depends on the concentration of hydrogen ion, it is possible to use this system to measure pH. A sample solution of unknown pH was saturated with quinohydrone and a platinum electrode was dipped into it. If the potential of the such electrode was found to be +0.22V with respect to SCE, what was the PH of the sample solution?
[Given : $\dfrac{2.303RT}{F}=0.06$]

\qn{2} 10 mL of H$_2$A (weak diprotic acid) solution is titrated against 0.1M NaOH. pH of the solution is plotted against volume of strong base added and following observation is made.
\qfig{p1_chem_s3_q02_1}
If pH of the solution at first equivalence point is pH$_1$ and at second equivalence point is pH$_2$.
Calculate the value of (pH$_2-$pH$_1$) at $25^\circ$C
Given: For H$_2$A, $pK_{a_1}=4.6$ and $pK_{a_2}=8$, $\log25=1.4$

\qn{3} Select the number of incorrect reaction(s).

(i) 2XeF$_6$ + 3SiO$_2$ $\longrightarrow$ 2XeO$_3$ + 3SiF$_4$

(ii) XeF$_4$ + Pt $\longrightarrow$ Xe + PtF$_4$

(iii) XeF$_4$ + (N(CH$_3$)$_4$)F $\longrightarrow$ N(CH$_3$)$_4$XeF$_5$

(iv) XeF$_2$ + AsF$_5$ $\longrightarrow$ (XeF)$^+$[AsF$_6$]$^-$

(v) 2XeF$_6$ + 3H$_2$O $\longrightarrow$ Xe + XeO$_3$ + 6HF + 3F$_2$

(vi) 2XeF$_6$ + 16NaOH $\longrightarrow$ Na$_4$XeO$_6$ + Xe + 8H$_2$O + 12NaF + O$_2$

(vii) XeF$_6$ + BF$_3$ $\longrightarrow$ (BF$_2$)$^+$(XeF$_7$)$^-$

\qn{4} Find the total number of geometrical isomers of octahedral complex.
[M(AB)(CD)(EF)]$^{n+}$ where (AB), (CD) and (EF) are unsymmetrical bidentate ligands.

\qn{5} How many of the following amino acids are considered as basic amino acids?

(i) Aspartic acid, (ii) Serine (iii) Tyrosine (iv) Cysteine

(v) Phenylalanine (vi) Histidine (vii) Valine

(viii) Arginine (ix) Alanine

\qn{6} How many of the following are prepared by condensation polymerization?

(i) Decron (ii) Cellulose (iii) Nylon-6,6 (iv) PVC (v) SBR (vi) PTFE

(vii) Bakelite (viii) Novolac (ix) Neoprene

\partsec{3}{MATHEMATICS}
\sect{Section-I: One or More Correct}

\qn{1} Consider the planes\\
P$_1$ : $x = y\sin\alpha + z\sin\beta$\\
P$_2$ : $y = z\sin\theta + x\sin\alpha$\\
P$_3$ : $z = x\sin\beta + y\sin\theta$\\
If P$_1$,P$_2$,P$_3$ intersect in a line $\dfrac{x}{\cos\theta}=\dfrac{y}{\cos\beta}=\dfrac{z}{\cos\alpha}$ then
\begin{qoptions}
\opt{A}{$\sin^2\alpha+\sin^2\beta+2\sin\theta\sin\alpha\sin\beta=\cos^2\theta$}
\opt{B}{$\sin^3\alpha+\sin^3\beta+\sin^3\theta=3\sin\alpha\sin\beta\sin\theta$}
\opt{C}{$\theta+\alpha+\beta=2n\pi+\dfrac{\pi}{2}$}
\opt{D}{$\theta+\alpha+\beta=\pi$}
\end{qoptions}

\qn{2} If f(x) is a polynomial of degree three with leading coefficient as unity such that f(1) = 1, f(3) = 9 and f$'$(1) = 2, then
\begin{qoptions}
\opt{A}{f$'$(x) = 2x has atleast one root between (1, 3)}
\opt{B}{f$''$(x) = 2 has atleast one root between (1, 7/3)}
\opt{C}{Product of roots of f(x) = 0 is 3.}
\opt{D}{Sum of roots of f(x) = 0 is 4.}
\end{qoptions}

\qn{3} Which of the following statement is correct ?
\begin{qoptions}
\opt{A}{$30! > 3!\,5!\,10!\,15!$}
\opt{B}{$4\sin^{-1}\left(\dfrac14\right)+3\sin^{-1}\left(\dfrac13\right)>7\sin^{-1}\left(\dfrac27\right)$}
\opt{C}{Consider a real system of equations $x^2+y^2=4$, $|x|+|y|=k$. Then the system has always even number of solutions.}
\opt{D}{$f(x)=\begin{cases}1-|x|, & |x|\le1\\ 0, & |x|>1\end{cases}$ and $g(x)=f(x-1)+f(x+1)$. Then the number of points where $g(x)$ is not differentiable is 4.}
\end{qoptions}

\qn{4} Consider the equation $Z^{200}-\lambda Z^{199}+1=0$. Let $Z_1,Z_2,Z_3,\dots,Z_{200}$ be non-zero complex numbers and $Z^{200}-\lambda Z^{199}+1=(Z-Z_1)(Z-Z_2)(Z-Z_3)\dots(Z-Z_{200})$ then
\begin{qoptions}
\opt{A}{Imaginary part of $\displaystyle\sum_{j=1}^{200}\sum_{k=1}^{200}\dfrac{Z_j}{Z_k}$ is 0}
\opt{B}{Imaginary part of $\displaystyle\sum_{j=1}^{200}\sum_{k=1}^{200}\dfrac{Z_j}{Z_k}$ is 1}
\opt{C}{Real part of $\displaystyle\sum_{j=1}^{200}\sum_{k=1}^{200}\dfrac{Z_j}{Z_k}$ is 0}
\opt{D}{Real part of $\displaystyle\sum_{j=1}^{200}\sum_{k=1}^{200}\dfrac{Z_j}{Z_k}$ is 1}
\end{qoptions}

\qn{5} If $(x-a)\cos\theta+y\sin\theta=(x-a)\cos\phi+y\sin\phi=a$ and $\tan(\theta/2)-\tan(\phi/2)=2b$, then
\begin{qoptions}
\opt{A}{$y^2=2ax-(1-b^2)x^2$}
\opt{B}{$\tan\dfrac{\theta}{2}=\dfrac1x(y+bx)$}
\opt{C}{$y^2=2bx-(1-a^2)x^2$}
\opt{D}{$\tan\dfrac{\phi}{2}=\dfrac1x(y-bx)$}
\end{qoptions}

\qn{6} Let $f(x)=\begin{cases}5x, & x\le2\\[2pt] \displaystyle\int_0^x(4+|2-t|)\,dt, & x>2\end{cases}$, then :
\begin{qoptions}
\opt{A}{$f(x)$ is continuous but not differentiable at $x=2$.}
\opt{B}{$f(x)$ is not continuous at $x=2$.}
\opt{C}{$f(x)$ is differentiable everywhere $\forall\ x\in(2,\infty)$}
\opt{D}{The right hand derivative of $f(x)$ at $x=3$ does not exist.}
\end{qoptions}

\sect{Section-II: Numerical Value}

\qn{1} The differential equation $\dfrac{dy}{dx}=\dfrac{(x+y-3)^2-(y-5)^2}{(x+2)^2}$ determines a curve $y=f(x)$ and $f(0)=1$, then $|f(3)|$ is equal to :

\qn{2} Let $f(x)=\dfrac{1}{e^x+8e^{-x}+4e^{-3x}}$, $g(x)=\dfrac{1}{e^{3x}+8e^{x}+4e^{-x}}$ and $\displaystyle\int(f(x)-2g(x))\,dx=h(x)+c$ (where `c' is constant of integration) and $\displaystyle\lim_{x\to\infty}h(x)=\dfrac{\pi}{4}$. If $h(0)=\dfrac1a\tan^{-1}\left(\dfrac bc\right)$ (where $a\in N$, $b$ and $c$ are coprime), then the value of $\left(\dfrac{a+c}{b}\right)$ is :

\qn{3} If $f(x)=\displaystyle\lim_{n\to\infty}\tan^{-1}\left(4n^2\left(1-\cos\dfrac xn\right)\right)$ and $g(x)=\displaystyle\lim_{n\to\infty}\dfrac{n^2}{2}\ell n\cos\left(\dfrac{2x}{n}\right)$, then $\displaystyle\lim_{x\to0}\dfrac{e^{-2g(x)}-e^{f(x)}}{x^6}$ equals to

\qn{4} Let A,B and C be 3 square matrices of order $3\times3$. If $A^3-6A^2+7A+B=0$ and $(\mathrm{adj}(\mathrm{adj}B))=C$, where $A=\begin{bmatrix}1&0&2\\0&2&1\\2&0&3\end{bmatrix}$, then $\left(\dfrac{\text{sum of elments of }C}{10}\right)$ is

\qn{5} Let the line $y-\sqrt3x+3=0$ cuts the parabola $2y^2=2x+3$ at A and B. If $P\left(\sqrt3,0\right)$, then value of $|PA-PB|$ is [where PA denotes distance between points P and A.]

\qn{6} The probability that Ramu's Gardner will forget to water the rose bush is $\dfrac23$. The rose bush is in questionable condition. Any how if watered, the probability of its withering is $\dfrac12$ and if not watered then the probability of its withering is $\dfrac34$. Ramu went out of station and after returning he finds that rose bush has withered. The probability that the Gardner did not water the rose bush is

\sect{Section-III: Single Digit Integer (0-9)}

\qn{1} Tangents are drawn from any point on straight line $y=x+8$ to the auxiliary circle of $x^2+16y^2=16$. If A and B are the points of contact of these tangent and P,Q are corresponding points of A and B on the ellipse respectively. If locus of mid point of PQ is $x^2+\alpha y^2+\beta x+\gamma y=0$ then value of $\dfrac{(\alpha+\beta-\gamma)}{13}$ is

\qn{2} If the sum of the series $\sec^2x+\dfrac{1}{2^2}\sec^2\left(\dfrac x2\right)+\dfrac{1}{2^4}\sec^2\left(\dfrac{x}{2^2}\right)+\dots+\dfrac{1}{2^{2n-2}}\sec^2\left(\dfrac{x}{2^{n-1}}\right)$ is equal to $a^2\mathrm{cosec}^2ax-\dfrac{4}{4^n}\mathrm{cosec}^2\dfrac{x}{2^{n-1}}$, $a\in N$, then the value of `a' is

\qn{3} Let $(x_0,y_0,z_0)$ be the solution of equation
\begin{align*}
3x-4y+\dfrac{1}{xy}&=12\\
4z-12x+\dfrac{1}{zx}&=3\\
12y-3z+\dfrac{1}{yz}&=4
\end{align*}
Then $\dfrac{z_0}{4x_0y_0}$ is

\qn{4} Let $g_1,g_2,g_3:(0,\infty)\to R$ be defined by $g_1(x)=\displaystyle\int_0^x\prod_{i=1}^{19}(t-i)^i\,dt$, $x>0$\\
and $g_2(x)=98(x-1)^{50}-100(x-1)^{49}+2450$, $x>0$\\
and $g_3(x)=(x-1)^{2020}(x-2)^{2021}(x-3)^{2022}$, $x>0$\\
Let $\alpha_i$ and $\beta_i$ respectively denote the number of points of local minima and the number of points of local maxima of the function $g_i$, $i=1,2,3$ in $(0,\infty)$ and if the value of $\alpha_1\beta_1+\alpha_2\beta_2+\alpha_3\beta_3$ is $\lambda$ then $\lambda-20$ equal to (where $\displaystyle\prod_{i=1}^{n}a_i=a_1a_2a_3\dots a_n$, $n\in N$)

\qn{5} Let $f:A\to A$ where $A=\{1,2,3,4,5,6,7\}$, then number of functions $f$ such that $f(f(f(x)))=x\ \forall\ x\in A$ is $\lambda$ then the value of $\dfrac{\lambda}{39}$ is

\qn{6} Find the number of elements contained in the range of the function $f(x)=\left[\dfrac x6\right]\left[\dfrac{-6}{x}\right]\ \forall\ x\in(0,30]$ (where [.] denotes greatest integer function)

\end{multicols}
\end{document}
